The quest for mastery of knowledge has been a foundational aspiration in the development of artificial intelligence systems. Historically, seminal works by \citet{McCarthy1963} and \citet{Newell1976} have underscored the significance of knowledge representation and reasoning in AI systems. For instance, the Cyc project embarked on an ambitious journey to codify common-sense knowledge, aiming to provide AI systems with a comprehensive understanding of the world \citep{Lenat1995}. Concurrently, endeavors like the WordNet project by \citet{Miller1990} sought to create lexical databases that capture semantic relationships between words, thereby aiding AI systems in grasping the nuances of human language. 

Amidst these pioneering efforts, the emergence of Large Language Models (LLMs), such as ChatGPT \citep{chatgpt}, GPT-4 \citep{GPT4} and LLaMA \citep{llama,llama2}, has been seen as a significant leap in both academics and industries, especially towards AI systems possessing vast factual knowledge \citep{LAMA,de-cao-etal-2021-editing,GPT4}. 
The advantages of using LLMs as knowledge bases carriers are manifold. Firstly, they reduce the overhead and costs associated with building and maintaining dedicated knowledge bases \citep{petroni2019language,alkhamissi2022review,KSSM}. Additionally, LLMs offer a more flexible approach to knowledge processing and utilization, allowing for context-aware reasoning and the ability to adapt to novel information or prompts \citep{sun2023beamsearchqa,huang2023reasoning}. 
Yet, with their unparalleled capabilities, concerns have arisen about the potential of LLMs to generate non-factual or misleading content \citep{Bender2021,GPT4,bubeck2023sparks}. In light of these advancements and challenges, this survey seeks to delve deeply into the LLMs, exploring both their potential and the concerns surrounding their factual accuracy.
%\wjdd{Can we draw a cute plot here to show what is factuality of LLMs?} {response: We will add one when submitting}

Understanding the factuality of Large Language Models is more than just a technical challenge; it's essential for the responsible use of these tools in our daily lives. As LLMs become more integrated into services like search engines \citep{BingChat}, chatbots \citep{chatgpt,bard}, and content generators \citep{ChatLaw}, the information they provide directly influences decisions, beliefs, and actions of millions of people. If an LLM provides incorrect or misleading information, it can lead to misunderstandings, spread false beliefs, or even cause harm, especially for those domains that demand high factual accuracy~\citep{ling2023beyond}, such as health ~\citep{thirunavukarasu2023large, tang-etal-2023-aligning}, law~\citep{huang2023lawyer}, and finance~\citep{BloombergGPT}. For instance, a physician relying on an LLM for medical guidance might inadvertently jeopardize patient health, a corporation leveraging LLM insights might make ill-informed market decisions, or an attorney misinformed by an LLM might falter in legal proceedings \citep{curran2023hallucination}.
In addition, with the advancement of LLM-based agents, the factuality of LLMs is becoming even more potent. A driver or an autonomous driving car might rely on LLM-based agents for planning or driving, where serious factual mistakes made by LLMs could cause irreversible damage.
By studying the factuality of LLMs, we aim to ensure that these models are both powerful and trustworthy.

A surge of research has been directed towards evaluating LLMs' factuality, which encompasses diverse tasks like factoid question answering and fact checking. Beyond evaluation, efforts to improve the factual knowledge of LLMs have been notable. Strategies have ranged from retrieving information from external knowledge bases to continual pretraining and supervised finetuning. Yet, despite these burgeoning efforts, a holistic overview that covers the full spectrum of factuality in LLMs remains elusive. 
While there are existing surveys in the field, such as those by \citet{chang2023survey} and \citet{aligning_llm_human}, that delve into the evaluation of LLMs and their factuality, they only scratch the surface of the broader landscape. There are also a bunch of recent studies focusing on hallucinations in LLMs \citep{rawte2023survey,zhang2023sirens,rawte2023survey,ye2023cognitive,Huang2023ASO}. But we differentiate between the hallucination issue and the factuality issues in Sec \ref{sec:definition}. Moreover, these surveys often overlook key areas we emphasize, like domain-specific factuality or the challenge of outdated information. While \citet{ling2023domain} explores domain specialization in LLMs, our survey takes a more expansive look at the broader issues of factuality.
To the best of our understanding, our work is the first comprehensive study on the factuality of large language models.
% \wjdd{Factuality relates to many aspects such as evaluation and value alignment. Maybe we should cite some existing survey papers here and discuss our difference. Then, we can move on to the next paragraph. R: I am not sure whether we should discuss the related survey here, or we shall choose another place to discuss them?}
% \citep{Hallucination_Survey,LLMSurvey}, \textbf{Outdated Information} \cite{WebGPT,WebCPM}, and \textbf{Domain-Specificity} (e.g., Health \citep{DoctorGLM,huatuo}, Law \citep{ChatLaw}, Finance \citep{BloombergGPT}).

This survey aims to offer an exhaustive overview of the factuality studies in LLMs, delving into four key dimensions: Sec 2) The definition and impact of the factuality issue \citep{Pranshu2023,nori2023capabilities}; Sec 3) Techniques for evaluating factuality and its quantitative assessment~\citep{C-Eval, Min2023-FActScore}; Sec 4) Analyzing the underlying mechanisms of factuality in LLMs and identifying the root causes of factual errors~~\citep{liu2023lost, kotha2023understanding}; and Sec 5) Approaches to enhance the factuality of LLMs~\citep{multiagent_debate, he2022rethinking}. Notably, we categorize the use of LLMs into two primary settings: LLMs without external knowledge, such as ChatGPT \citep{chatgpt} and Retrieval-Augmented LLMs, such as BingChat \citep{BingChat}. The complete structure of this survey is illustrated in Figure \ref{fig:tree}.
Through a detailed examination of existing research, we seek to shed light on this critical aspect of LLMs, helping researchers, developers, and users harness the power of these models responsibly and effectively.